% Motivation and learning outcomes.
% The outcomes are the outcomes of Day 7 in the syllabus.
\begin{frame}{Why this day matters}
  \begin{itemize}
    \item Days 1 to 6 made the model smaller. From today the device is the
          Raspberry Pi 5: a Linux computer with four processor cores and 8 GB
          of RAM.
    \item On this device, the same model can run in more than one runtime. The
          runtime, its kernels, the number of threads, and the precision
          change the latency.
    \item A faster processor does not help a workload that waits for memory.
          The Roofline model shows which of the two limits applies.
    \item In the lab you run MobileNetV2 on the Raspberry Pi, export one model
          to ONNX and to LiteRT, and measure the latency for each runtime,
          thread count, and precision.
  \end{itemize}
\end{frame}

\begin{frame}{Learning outcomes}
  After this day you can:
  \begin{itemize}
    \item Explain how SIMD, accelerators, and memory bandwidth limit inference
          speed.
    \item Name the main graph optimizations and state what each one does.
    \item Export a PyTorch model to ONNX and to LiteRT.
    \item Compare inference runtimes on the Raspberry Pi.
    \item Explain the trade-off between static and dynamic input shapes.
  \end{itemize}
\end{frame}
